\chapter*{导\quad 论}\phantomsection\addcontentsline{toc}{chapter}{导\quad 论}\markboth{导\quad 论}{导\quad 论}这本书要讲的是一座建在月球背面地下、运转了两万余年的城市，以及它怎样在公元前二二一年把自己清空。动手翻阅之前，有件事必须放在头里讲清：笔者没有去过月面。桌面上的全部材料，不过是二十九帧早期轨道影像、一根湖底岩芯、几百片刻着小篆的薄合金板、四十六枚折断的钻头，还有几份编目年份横跨上世纪六十年代到八十年代的测绘简报。这本书能站到什么位置，全看这些冷物什自己肯开口说出多少话。

按丛书原来的分工，关于黄土高原上的着陆残片以及中土王朝后世的几度震荡，已经另有《战争的雷霆》一书站在地面一侧作过梳理；关于月球正面静海一隅出土的三十七片残简，也已有《月之都——〈广寒宫残简〉辑佚与考释》专册校注。本书的体裁是一部覆盖全月面的概论，重心放在正面与背面连成一体的整座城池、城内维持两万年的制度、向地面投放历法的社会实验，以及撤离前后那些散落到中原与日本列岛的人事与名号。即便读者此前没有读过上面提到的任何一本，依然可以从第一章第一帧照片读起，不必在手边另备前书对照。

\section*{一、材料层级与论证方法}本书把所有依据硬划成两层。第一层是一手实物与原生档册：轨道摄影与重力场测量、贝加尔湖沉积岩芯、格陵兰冰芯参考、合金板《章程总目》《工籍残卷》《规程残卷》《隔离舱管理档》《撤离签符与登船核对册》、净秽山北侧浅层岩芯与断钻头、有穷氏故地地下空腔测绘，以及东郡出土的高温熔蚀残片。全书十章的骨架全部压在这一层材料上。第二层是传世文献与后世神话，包括《尚书》《诗经》《左传》《天问》《山海经》《淮南子》《史记》《古事记》《日本书纪》与《竹取物语》。凡涉及第二层材料之处，一律只拿来跟一手证据对齿或者充作旁证，绝不用神话说事来代替实物。

在取材与断代上，本书守一条简单的规矩：只要传世文本确有其书、传流有序，便按它自身的史料性质定位，不因为成书晚就一笔否定。《竹取物语》定稿迟至九、十世纪，《古事记》与《日本书纪》也晚在八世纪初才写成文字，距离公元前二一〇年东渡隔了九百余年；它们记不下当时的月份与人数，却留住了前文那批名字在口传里变形后的最后落点。真正该剔除的，是后人凭空补出来的寓意；晚出文本本身的史料层位，不能因为定稿年代偏晚就随手抹去。

为了避免只靠几段孤立铭文就编出一整部通史，全书采取一章认一件主证物的做法。前五章由外往里剥开这座城池：第一章从背面轮廓与正面一角入手，量出人工边界与壳层结构；第二章拿湖底岩芯与题头接出两万一千三百年前的迁移上限；第三章读透四百一十二枚总目合金板，查清为什么两万年没有修改过一条律令；第四章与第五章分别清理《名籍》《工籍》两套对不上的账目和《净秽》卷一百零四条，指明被不断抽离出城的「变化」到底指什么。第六章起转入城外与地面：先查净秽山隆起与蓬莱药的由来，第七章与第八章顺着十一张《规程残卷》和一套旧历抄本，还原公元前二〇〇〇年前后的历法实验与纯化事件；第九章与第十章则翻开隔离舱一千七百六十四年的记录与撤离名册，接上公元前二二一年的大撤退、前二一〇年的跨海东渡，以及那些名号在随后两千年里的层累与消散。

读者在翻阅时大抵会觉得前半部分核得极细、后半部分步子偏快。这种前紧后松的节律实由手头材料的自然分布决定：前五章面对的是保留了尺寸、刻痕与地质剖面的硬器物，后五章牵涉到的大多是残缺的规程、断角的金属板以及只剩一两行摘要的传世记载。每章收尾都设有一节专门盘点材料缺口的文字，凡是没有实物托底、仅靠上下文接榫之处，笔者一律在各章之内直接标出这一步属于猜测，绝不用模棱两可的修辞替空白代劳。

\section*{二、本书的边界}概论最怕越界去抢专论的饭碗，因此全书定下三条清晰的边界。这三条边界把静海七号细部、中原王朝正史以及消忘现象的物理机制各自留给同系专著，只取与全月面主线咬合最紧的环节入书。

第一，凡是静海七号那三十七片残简的文字校勘、残画考订与城内水风管网的细部复原，原则上归《月之都——〈广寒宫残简〉辑佚与考释》承担。本书凡引及这批残简，只交代其大意与编号，并在注中注明出处于彼书，不再重抄整段释文与考据过程。第二，中国地上王朝在公元前二二一年之后经历的各轮技术承替、次生黄土沉积及政治震荡，主体已经在《战争的雷霆》中处理完毕，本书只在中原降落人数、旧址接管与「华夏神族」名号攀附这几处必要的接口上落墨，不另起炉灶重讲一遍中国史。

第三，全书会反复碰到一种名为「记忆消除场」的现象：为什么六十年代拍得清清楚楚的月面建筑底片在档案库里睡了几十年无人过问，为什么隔离舱里的服药者不断失忆，又为什么挂了两千年的名号在清代以后两三代之内就失了传。对于这些现象背后的物理与认知机理，本书一概不做解释。场的作用原理属于索科洛夫《记忆消除场的文献学后果》的研究范围，本书只在档案与文献留出空白的地方指一下手，把那些可见的后果记下来，其余留给专著去答。

全书自此以下进入正文，不依赖任何预设前提，也不预先替书名里的几个字作解。每一章都在前文已经坐实或标明存疑的地基上往前递进半步，走到哪里算到哪里，接不上的地方就留在那里。
